\section{Coordination hypothesis and safety conditions}
\label{sec:hypothesis}

Our proposal addresses a conditional question: if cross-organisational work becomes agent-mediated, what infrastructure would allow changing participants to discover one another, interpret prior exchanges and inspect the evidence for a decision? We do not predict that this shift will occur. Murmur connects portable identity, published schemas and attributable conversation records to support such work. Its single-operator coordination prototype demonstrates a path for exchanging and reconstructing messages, but cannot establish the benefit across independent organisations. We state the proposed benefit as a testable hypothesis.

\begin{quote}
\textbf{H1 (coordination benefit).} For agent-mediated work across independently operated organisations, a shared, attributable interaction record, supported by verifiable identities and published schemas, improves the completeness and timeliness of independent decision reconstruction and recovery under changing participation and interrupted communication, compared with point-to-point exchanges using equivalent authority, resource and retention controls.
\end{quote}

Here, \emph{shared} means that authorised participants can obtain and verify the relevant records under agreed policies. It does not require one globally ordered log or identical access to private information. \emph{Independent reconstruction} means establishing the accepted terms, their declared dependencies, the required authority evidence and the observed terminal decision without relying on the initiating agent's surviving process state. Recovery means reconciling an interrupted operation with its durable outcome; replaying the conversation alone does not recover a lost resource ledger.

The proposed mechanism combines discoverable participants and interpretable record types with explicit dependency preservation. Each request, delegation and response identifies the records on which it depends. A joining participant can retrieve those dependencies, and a recovering service can verify a retrieved outcome. Failure to retrieve an outcome does not establish that no decision was made. For commerce, the record must also identify the gateway evidence associated with a reservation or terminal decision. Published references make these relationships inspectable, while retention and access policies determine whether the evidence remains available.

The same mechanism creates a disclosure constraint: making evidence available can expose the private information used to produce it. Public and private exchanges therefore cannot be assumed to form one universally readable thread. Cross-boundary correlation remains private where exposing the link would reveal protected context, and release requires enforcement beyond an agent's judgement. A gain in reconstruction achieved by disclosing information outside the agreed policy would not support the proposed benefit. Murmur specifies these protections as design requirements; they are not implemented guarantees.

This explanation has a credible alternative. Signed point-to-point messages, a shared directory and durable local histories may provide the same reconstruction capability. Similarly, improved recovery may result entirely from the outbox and reservation rules rather than from AT Protocol publication. Comparisons must therefore retain those controls in both designs. AT Protocol is the candidate substrate examined here, not a necessary condition for the hypothesis.

The hypothesis predicts an advantage where participants join late, evidence originates in several organisations or a service restarts after a decision. It predicts less benefit for stable workflows operated by one trusted team, where a database and queue may provide the required context more simply. It would be weakened if equivalent point-to-point evidence gives the same reconstruction coverage and recovery time. A practical adoption claim would also fail if any improvement requires unacceptable replication cost, disclosure or operational complexity. These conditions separate a plausible coordination mechanism from a claim that federation is universally preferable.

The comparative claim in H1 is empirical. It cannot be proved from the transaction rules alone. We therefore distinguish it from three conditional safety properties of Murmur's initial-commitment protocol: budget preservation, exclusive terminal decisions and retained backing during delayed communication (Propositions~\ref{prop:budget-safety}--\ref{prop:reservation-backing}). Their proofs concern the specified transitions under trusted, durable enforcement. They establish neither correctness of a deployed implementation nor superiority of the coordination substrate. A comparator must satisfy the same safety conditions before its reconstruction and recovery performance is compared with Murmur.
